\documentclass[11pt]{article}
\usepackage{fontspec,amsmath,amssymb,amsthm,hyperref}
\usepackage[margin=1in]{geometry}
\setlength{\emergencystretch}{3em}
\title{Actual Brownian paths under a kinetic clock change}
\author{LangevinLean mathematical exposition}
\date{}
\begin{document}\maketitle
\noindent Source module: \texttt{KineticLangevin/BrownianScaling.lean}. Compile this proof document with XeLaTeX or LuaLaTeX. The accompanying HTML pairs the mathematical argument with the exact Lean statements and records the scope of the formalization.


Time change is a probability statement as well as a deterministic substitution in the drift. If \(c>0\), the process \(\widehat B_t=\sqrt c B_{t/c}\) has the same Brownian finite-dimensional laws, because its covariance is \(c\min(s/c,t/c)=\min(s,t)\). This module proves that assertion on the already constructed canonical source, retains independence of entire coordinate paths, and calculates the physical noise amplitude.


\section{Change the actual coordinate process and prove its Brownian law}


\texttt{clockBrownian} applies the clock directly to each coordinate of the canonical source. Mathlib’s proved Brownian scaling theorem uses \(a^{-1/2}B_{at}\); setting \(a=c^{-1}\) gives the desired formula. Lean simplifies the inverse square root and the double inverse to identify this process exactly. \texttt{IsPreBrownianReal} supplies all finite-dimensional Gaussian distributions and independent increments, rather than a single matching variance.



\noindent Lean declarations: \texttt{clockBrownian}, \texttt{clockBrownian\_isPreBrownian}.

\section{Preserve independence and recover the kinetic temperature}


Each entire changed-clock coordinate path is a measurable function of its own original coordinate path. Independence is therefore preserved by measurable composition. Finally, velocity scaling by \(c^{-1}\) changes the standard thermal input to \(c^{-1}\sqrt{2c\gamma}\,B_{t/c}=\sqrt{2\gamma/c^2}\,\widehat B_t\). The scalar proof uses the square-root product and quotient identities with their nonnegativity hypotheses. For \(c=\sqrt N\), this is exactly the physical particle temperature \(1/N\).



\noindent Lean declarations: \texttt{clockBrownian\_paths\_independent}, \texttt{kinetic\_clock\_noise\_amplitude}.

\end{document}
